\section{Introduccion del capitulo}
% =====================================================================

Este capitulo documenta el programa experimental construido sobre el
proyecto \texttt{MetaCARP} para evaluar empiricamente la calidad de
distintas estrategias de seleccion de operadores, mecanismos de
intensificacion, esquemas de penalizacion de la infactibilidad y
heuristicas de construccion inicial en el contexto del
\emph{Capacitated Arc Routing Problem} (CARP). El programa se
desarrollo durante el mes de mayo de 2026 y comprende siete
experimentos principales del grid completo $23\text{ instancias} \times
5\text{ repeticiones} \times 5\text{ metaheuristicas}$ (\textbf{$575$
corridas por experimento}), ademas de un conjunto de experimentos de
control que se discuten en la Seccion~\ref{sec:re-greedy}.

\subsection{Estructura del capitulo}

El capitulo sigue una organizacion \textbf{cronologica}, en el orden en
que los experimentos se planearon y ejecutaron. Esta eleccion no es
casual: cada experimento se origino como respuesta a una observacion
puntual del experimento anterior. Mantener la narrativa cronologica
permite reproducir el razonamiento que justifico cada hiperparametro,
cada cambio de selector y cada mecanismo nuevo. En particular, la
narrativa cronologica es la unica forma honesta de presentar el
hallazgo metodologico mas importante del trabajo, que se documenta en
la Seccion~\ref{sec:evaluador-canonico}: la deteccion de un artefacto
en el evaluador de costo de las metaheuristicas que inflaba
artificialmente todos los \emph{gap} reportados.

Las secciones se organizan como sigue:

\begin{itemize}
\item Las Secciones~\ref{sec:marco-experimental}--\ref{sec:operadores}
  describen el marco experimental: instancias, metricas,
  infraestructura, los cinco algoritmos base (SA, TS Simple, RTS, ABC,
  Cuckoo) y los operadores de vecindad.
\item La Seccion~\ref{sec:arquitectura-software} documenta la
  \emph{arquitectura de monkey-patching por capas}, una decision
  metodologica que permitio incrementar componentes sin tocar el
  codigo base de las metaheuristicas.
\item Las Secciones~\ref{sec:exp-strict}--\ref{sec:exp-budget}
  presentan los seis experimentos del primer ciclo, en orden
  cronologico: \texttt{strict\_intra\_inter},
  \texttt{lambda\_grid}, \texttt{aos\_pm}, \texttt{path\_relinking},
  \texttt{p\_inter\_pr} y \texttt{budget}.
\item La Seccion~\ref{sec:evaluador-canonico} expone el hallazgo
  critico: el evaluador canonico de las metaheuristicas forzaba
  orientacion fija $u \to v$ en cada arco, lo que producia
  mediciones de costo fisicamente incorrectas y por tanto \emph{gap}
  inflados en TODOS los experimentos previos.
\item La Seccion~\ref{sec:warmstart-greedy} introduce los dos
  componentes que cierran el gap: Path-Scanning como heuristica
  constructiva inicial y el evaluador con orientacion greedy.
\item La Seccion~\ref{sec:re-greedy} presenta los experimentos de
  control disenados para aislar el efecto del evaluador.
\item La Seccion~\ref{sec:conclusiones} consolida los hallazgos
  metodologicos y tecnicos del programa experimental.
\end{itemize}

\subsection{Aprendizajes que el capitulo busca documentar}

Los siguientes aprendizajes son el resultado neto del programa
experimental. Cada uno se justifica empiricamente en la seccion
correspondiente.

\begin{enumerate}
\item El \textbf{selector binario estricto} (factible $\to$ intra,
  infactible $\to$ inter) inhibe la exploracion: las metaheuristicas
  basadas en trayectoria (TS Simple, RTS) terminan con tasas de
  factibilidad final del 0--10\%.
\item Un \textbf{selector probabilistico moderado} ($p_{\text{inter}}
  = 0.20$ en estado factible) supera al binario estricto y a un
  selector adaptativo de tipo Probability Matching en las cinco MHs.
\item \textbf{Path Relinking truncado} contribuye especialmente en
  metaheuristicas basadas en poblacion (ABC, Cuckoo), donde la
  poblacion provee guias diversas; en metaheuristicas de trayectoria
  la contribucion es marginal o nula.
\item \textbf{Aumentar el presupuesto computacional} en un factor
  $\times25$ produce ganancias marginales (0--2\,pp); el cuello de
  botella en el regimen estudiado NO es el numero de iteraciones.
\item El verdadero cuello de botella fue un \textbf{artefacto en el
  evaluador}: la forzar orientacion canonica $u \to v$ por tarea
  inflaba el costo de cualquier solucion entre 10 y 50\,pp.
\item Una \textbf{construccion inicial heuristica} (Path-Scanning
  clasico de Golden, DeArmon y Baker~\cite{golden1983}) combinada con la
  correccion del evaluador permite alcanzar el BKS exacto en multiples
  instancias del benchmark.
\end{enumerate}

% =====================================================================
\section{Marco experimental}
\label{sec:marco-experimental}
% =====================================================================

\subsection{Conjunto de instancias}
\label{sec:instancias}

Las pruebas se realizaron sobre \textbf{23 instancias} de los
benchmarks GDB~\cite{golden1983} y KSHS~\cite{kiuchi1995}, ampliamente
utilizados en la literatura de CARP~\cite{golden1981}. La
Tabla~\ref{tab:instancias} resume sus caracteristicas.

\begin{table}[h]
\centering
\caption{Instancias del benchmark utilizadas.}
\label{tab:instancias}
\small
\begin{tabular}{lrrrrr}
\toprule
Instancia & $|V|$ & $|E_R|$ & $|V|$eh & $Q$ & BKS \\
\midrule
gdb1  & 12 & 22 & 5 &   5 &  316 \\
gdb2  & 12 & 26 & 6 &   5 &  339 \\
gdb3  & 12 & 22 & 5 &   5 &  275 \\
gdb4  & 11 & 19 & 4 &   5 &  287 \\
gdb5  & 13 & 26 & 6 &   5 &  377 \\
gdb6  & 12 & 22 & 5 &   5 &  298 \\
gdb7  & 12 & 22 & 5 &   5 &  325 \\
gdb10 & 12 & 25 & 4 &  10 &  275 \\
gdb12 & 9  & 19 & 7 &  20 &  458 \\
gdb13 & 7  & 35 & 6 &  27 &  536 \\
gdb14 & 7  & 21 & 5 &  21 &  100 \\
gdb15 & 7  & 21 & 4 &  37 &   58 \\
gdb16 & 8  & 28 & 5 &  24 &  127 \\
gdb17 & 8  & 28 & 5 &  41 &   91 \\
gdb19 & 8  & 11 & 3 &  27 &   55 \\
gdb20 & 11 & 22 & 4 &  27 &  121 \\
gdb21 & 11 & 33 & 6 &  27 &  156 \\
kshs1 & 8  & 15 & 3 & 110 & 14661 \\
kshs2 & 10 & 15 & 3 & 140 &  9863 \\
kshs3 & 6  & 15 & 3 & 110 &  9320 \\
kshs4 & 8  & 15 & 3 & 130 & 11498 \\
kshs5 & 8  & 15 & 3 &  75 & 10957 \\
kshs6 & 9  & 15 & 3 & 100 & 10197 \\
\bottomrule
\end{tabular}
\end{table}

En la tabla, $|V|$ es el numero de vertices, $|E_R|$ el numero de aristas
requeridas (las tareas), $|V|$eh el numero de vehiculos disponibles, $Q$
la capacidad de cada vehiculo y BKS (\emph{Best Known Solution}) el
mejor valor conocido en la literatura para esa instancia. Estas
instancias representan un rango amplio de tamanos y proporciones
$|E_R|/|V|$, lo que permite observar comportamientos diferenciados de
las metaheuristicas.

\subsection{Metrica de evaluacion}
\label{sec:metrica}

La metrica principal es el \emph{gap relativo al BKS}:

\begin{equation}
\gap(s) = \frac{c(s) - \BKS}{\BKS} \cdot 100\%
\label{eq:gap}
\end{equation}

donde $c(s)$ es el costo total de la solucion $s$ entregada por la
metaheuristica y BKS es el mejor valor conocido. Un gap del 0\%
significa que se igualo el BKS. Cuando se reporta el gap de un
experimento, sin mas calificadores, se entiende como la media
aritmetica del gap sobre las 23 instancias y las 5 repeticiones
ejecutadas.

Cuando se busca aislar la calidad maxima alcanzable se reporta el
\emph{best-of-N}: el minimo del gap entre las $N$ repeticiones de cada
instancia, promediado sobre las 23 instancias. Esta es la metrica de
referencia para concursos y rankings.

\subsection{Funcion de costo penalizada}
\label{sec:costo-penalizado}

Las metaheuristicas trabajan sobre el siguiente \emph{objetivo
penalizado}, que combina el costo puro con un termino de violacion de
capacidad:

\begin{equation}
\Phi_\lambda(s) = c(s) + \lambda \cdot \viol(s),
\label{eq:objetivo}
\end{equation}

donde $\viol(s) = \sum_{r \in s} \max(0, d(r) - Q)$ es el exceso total
de demanda sobre la capacidad $Q$ sumado sobre todas las rutas $r$ de
$s$, y $\lambda \ge 0$ es un factor de penalizacion que el experimento
\texttt{lambda\_grid} (Seccion~\ref{sec:exp-lambda}) calibra.

El valor $\lambda$ por defecto, instance-aware, es

\begin{equation}
\lambda_{\text{def}} = \max\bigl(10, \, 10 \cdot \mathrm{mediana}(D_{ij})\bigr),
\label{eq:lambda-def}
\end{equation}

donde $D_{ij}$ son las entradas positivas de la matriz de distancias
minimas de la instancia. Este valor escala con el tamano de la
instancia y garantiza que la penalizacion sea ni dominante ni
inocua.

\subsection{Infraestructura computacional}
\label{sec:infra}

Las corridas se ejecutaron en la siguiente plataforma:
\begin{itemize}
\item CPU: AMD Ryzen / Intel x86\_64 (segun nodo).
\item Sistema: Linux con kernel 6.17.
\item Python 3.12 con NumPy 1.26, NetworkX 3.4.
\item GPU: deshabilitada explicitamente (\texttt{usar\_gpu=False}) en
  todos los experimentos por incompatibilidad de \texttt{libnvrtc.so.12}.
\item Paralelismo: \texttt{ProcessPoolExecutor} con un worker por
  corrida, tantos workers como nucleos disponibles.
\item Tiempo estimado por experimento del grid completo: 30--90
  minutos en paralelo con 8--16 workers.
\end{itemize}

\subsection{Replicabilidad y semillas aleatorias}
\label{sec:semillas}

Las cinco metaheuristicas dependen de un generador \texttt{random.Random}
de Python. Todas las repeticiones invocan al wrapper con
\texttt{semilla=None}, lo que delega la inicializacion a
\texttt{os.urandom} (entropia del sistema operativo). Cada una de las
$575$ corridas por experimento utiliza por tanto una trayectoria
estocastica \emph{independiente}, pero NO reproducible: el campo
\texttt{semilla} del CSV almacena \texttt{None}, no la semilla
efectivamente sorteada por el SO.

La ausencia de \texttt{numpy.random} o \texttt{cupy.random} en el
codigo de las metaheuristicas garantiza que NO hay estado global
contaminante entre repeticiones. Esto se verifico inspeccionando con
\texttt{grep} todos los modulos del paquete \texttt{metacarp}.

% =====================================================================
\section{Algoritmos base}
\label{sec:algoritmos-base}
% =====================================================================

El programa experimental utiliza cinco metaheuristicas, divididas en dos
familias: metaheuristicas de \emph{trayectoria} (operan sobre una unica
solucion actual) y metaheuristicas \emph{poblacionales} (mantienen un
conjunto de soluciones simultaneamente).

\subsection{Recocido Simulado (SA)}

SA~\cite{kirkpatrick1983} es una metaheuristica de trayectoria que acepta
movimientos deteriorantes con probabilidad $\exp(-\Delta/T)$, donde $T$ es la
temperatura que decrece en cada nivel segun
$T_{k+1} = \alpha \cdot T_k$. El esquema implementado anade
\emph{reheating} (cuando se acumulan \texttt{patience} niveles sin
mejora, $T$ se reinicia a $\text{reheat\_factor} \cdot T_{\text{init}}$).

\subsection{Busqueda Tabu Simple}

TS Simple~\cite{glover1986, glover1997} mantiene una lista \emph{tabu} de movimientos prohibidos
durante un \texttt{tabu\_tenure} iteraciones. En cada iteracion explora
un vecindario de tamano \texttt{tam\_vecindario} y elige el mejor vecino
no tabu (o que cumpla la condicion de aspiracion: mejorar la mejor
solucion global).

\subsection{Busqueda Tabu Reactiva (RTS)}

RTS~\cite{battiti1994} extiende TS Simple con dos mecanismos adicionales:
(i) \emph{tenure adaptativo}: aumenta el \texttt{tabu\_tenure} cuando
detecta ciclos, lo reduce cuando hay diversificacion;
(ii) \emph{escape}: si una solucion se visita mas de
\texttt{umbral\_repeticiones\_escape} veces, se aplica una perturbacion
de \texttt{num\_movimientos\_escape} pasos aleatorios.

\subsection{Colonia Artificial de Abejas (ABC)}

ABC~\cite{karaboga2005} es una metaheuristica poblacional. Mantiene un conjunto de
\texttt{num\_fuentes} soluciones (las \emph{fuentes de alimento}). En
cada ciclo, cada abeja \emph{empleada} explora un vecino de su fuente;
las abejas \emph{observadoras} eligen una fuente con probabilidad
proporcional a su fitness y la perturban; cuando una fuente no mejora
en \texttt{limite\_abandono} ciclos, se abandona y se sustituye por una
solucion aleatoria nueva (\emph{scout}).

\subsection{Cuckoo Search}

Cuckoo Search~\cite{yang2009} mantiene un conjunto de \texttt{num\_nidos} soluciones.
En cada iteracion, cada nido genera un nuevo huevo aplicando un
\emph{Levy flight} (\texttt{pasos\_levy\_base} pasos aleatorios con
distribucion de cola pesada). El nuevo huevo reemplaza al nido si
mejora la calidad. Una fraccion \texttt{pa\_abandono} de los peores
nidos se abandonan al final de cada iteracion y se sustituyen.

Los pseudocodigos detallados de las cinco MHs se encuentran en el
directorio \texttt{pseudocodigos/} del repositorio del proyecto y son
identicos a los presentados en el Reporte previo
\emph{Seleccion de operadores y su efecto}; este capitulo se concentra
en las decisiones experimentales que se construyeron \emph{encima} de
estas implementaciones.

% =====================================================================
\section{Operadores de vecindad}
\label{sec:operadores}
% =====================================================================

CARP se modela sobre rutas: una solucion es una lista de rutas, donde
cada ruta es una secuencia de tareas servidas $(t_1, t_2, \dots, t_k)$.
Los \emph{operadores de vecindad} producen perturbaciones locales.

\subsection{Operadores intra-ruta}

Modifican el orden de las tareas \emph{dentro de una unica ruta}; no
cambian la asignacion de tareas a vehiculos, por lo que NUNCA introducen
violaciones de capacidad. Tres operadores intra estan implementados:

\begin{itemize}
\item \texttt{relocate\_intra}: retira una tarea de una posicion y la
  inserta en otra dentro de la misma ruta.
\item \texttt{swap\_intra}: intercambia dos tareas no consecutivas de
  la misma ruta.
\item \texttt{2opt\_intra}: invierte un segmento de la ruta entre dos
  posiciones $i$ y $j$.
\end{itemize}

\subsection{Operadores inter-ruta}

Mueven tareas entre rutas distintas. Pueden cambiar las cargas de las
rutas, por lo que SI pueden introducir o reparar violaciones de
capacidad. Seis operadores inter estan implementados, de los cuales se
usaron cinco como el conjunto \texttt{OPERADORES\_STRICT\_5} en los
experimentos del programa:

\begin{itemize}
\item \texttt{relocate\_inter}: retira una tarea de una ruta y la
  inserta en otra. (\emph{No incluido en el conjunto STRICT\_5}.)
\item \texttt{swap\_inter}: intercambia una tarea de la ruta $r_1$ con
  una de la ruta $r_2$.
\item \texttt{2opt\_star}: dada una posicion $i$ en $r_1$ y $j$ en $r_2$,
  reemplaza ambos sufijos por sus contrapartes cruzadas. (\emph{Incluido en STRICT\_5}.)
\item \texttt{cross\_exchange}: intercambia dos segmentos entre dos rutas.
\item \texttt{or\_opt\_2}: mueve un segmento de longitud 2 entre rutas.
\item \texttt{or\_opt\_3}: mueve un segmento de longitud 3 entre rutas.
\end{itemize}

El conjunto \texttt{OPERADORES\_STRICT\_5} usado en todos los
experimentos del programa esta compuesto por
$\{\texttt{relocate\_intra}, \texttt{swap\_intra},
\texttt{swap\_inter}, \texttt{2opt\_star}, \texttt{cross\_exchange}\}$:
dos intra y tres inter. Esta seleccion se fundamenta en analisis
previos de tasa de mejora por operador (Seccion 12 del notebook del
proyecto), no se repite aqui.

\subsection{Selector de grupo de operadores}
\label{sec:selector}

Cada iteracion de cualquier metaheuristica requiere decidir si el
proximo vecino se generara con un operador intra o uno inter, y luego
elegir el operador concreto dentro del grupo. El \emph{selector}
gobierna esa primera decision.

El selector original (\texttt{seleccionar\_grupo\_operadores\_inter\_intra}
en \texttt{metacarp.metaheuristicas\_utils}) usa dos parametros:

\begin{itemize}
\item $\alpha_{\text{inter}}$: probabilidad de elegir inter cuando la
  solucion actual VIOLA capacidad.
\item $p_{\text{inter}}$: probabilidad de elegir inter cuando la
  solucion actual es FACTIBLE.
\end{itemize}

El experimento \texttt{strict\_intra\_inter\_20260524} sustituye este
selector probabilistico por uno binario estricto; el experimento
\texttt{p\_inter\_pr\_20260528} regresa a uno probabilistico con
$p_{\text{inter}} = 0.20$ moderado. Estas decisiones se justifican en
sus respectivas secciones.

% =====================================================================
\section{Arquitectura de monkey-patching por capas}
\label{sec:arquitectura-software}
% =====================================================================

Una decision metodologica importante del programa experimental fue
NUNCA modificar el codigo de las metaheuristicas base
(\texttt{metacarp/recocido\_simulado.py}, etc.) durante el desarrollo
de las variantes experimentales. En su lugar, cada experimento
implementa sus cambios via \emph{monkey-patching} en tiempo de
ejecucion, aplicado al inicio de cada worker.

\subsection{Motivacion}

El monkey-patching ofrece tres ventajas en el contexto experimental:

\begin{enumerate}
\item \textbf{Comparabilidad}: cada experimento corre sobre exactamente
  el mismo codigo base, garantizando que los efectos observados son
  imputables a la modificacion del experimento.
\item \textbf{Composabilidad}: capas distintas (selector, kick reactivo,
  Path Relinking, evaluador) se pueden combinar sin reescribir el
  codigo de las metaheuristicas.
\item \textbf{Trazabilidad}: cada modulo experimental contiene en una
  funcion \texttt{aplicar\_patch\_X(modulo\_mh)} una documentacion
  ejecutable de exactamente que simbolos modifica.
\end{enumerate}

\subsection{Estructura general de un experimento}

Cada experimento del programa expone un par
\texttt{(metacarp/X\_YYYYMMDD.py, scripts/run\_<mh>\_X\_YYYYMMDD.py)}:

\begin{itemize}
\item El modulo en \texttt{metacarp/} define las funciones de
  reemplazo y una funcion \texttt{aplicar\_patch\_X(modulo\_mh)} que
  sustituye los simbolos relevantes en el modulo de la MH.
\item El script en \texttt{scripts/} implementa el bucle del grid
  experimental: itera sobre las $23 \times 5$ tareas, las distribuye
  con \texttt{ProcessPoolExecutor}, y dentro de cada worker llama
  \texttt{aplicar\_patch\_X(...)} ANTES de importar el wrapper
  \texttt{X\_desde\_instancia} de la metaheuristica.
\end{itemize}

\subsection{Composicion de capas}

Los experimentos finales del programa (\texttt{p\_inter\_pr\_20260528},
\texttt{budget\_20260528}, \texttt{warmstart\_greedy\_20260529}) APILAN
multiples capas de monkey-patch. La Figura~\ref{fig:capas} esquematiza
la composicion (formato pseudo-ascii):

\begin{verbatim}
+----------------------------------------------+
| Capa 5: Evaluador greedy (orient. dinamica)  |  <- warmstart_greedy
+----------------------------------------------+
| Capa 4: Path-Scanning warmstart              |  <- warmstart_greedy
+----------------------------------------------+
| Capa 3: Path Relinking truncado (p_pr=0.5)   |  <- pr, p_inter_pr, budget
+----------------------------------------------+
| Capa 2: Selector p_inter probabilistico      |  <- p_inter_pr, budget
|         (sobrescribe selector strict)         |
+----------------------------------------------+
| Capa 1: Selector binario estricto             |  <- strict, pr
|         + kick reactivo                       |
+----------------------------------------------+
|     Base MH (SA, TS, RTS, ABC, Cuckoo)        |
+----------------------------------------------+
\end{verbatim}

El \textbf{orden de aplicacion es critico} en dos puntos:

\begin{itemize}
\item La Capa 2 (\texttt{p\_inter}) sobrescribe el selector binario
  que la Capa 3 (PR) deja instalado. Por tanto, el orden correcto es:
  primero PR, luego \texttt{p\_inter}, para que el selector
  probabilistico prevalezca.
\item Las Capas 4 y 5 (Path-Scanning y evaluador greedy) son
  ortogonales al stack inferior y se pueden aplicar antes o despues,
  pero deben aplicarse en cada worker antes de la importacion del
  wrapper \texttt{<mh>\_desde\_instancia}.
\end{itemize}

% =====================================================================
\section{Experimento 1: strict\_intra\_inter\_20260524}
\label{sec:exp-strict}
% =====================================================================

\subsection{Hipotesis y motivacion}

El primer experimento del programa parte de la siguiente observacion:
el selector probabilistico original
($\alpha_{\text{inter}} = 0.80$, $p_{\text{inter}} = 0.60$) puede
proponer un operador inter incluso cuando la solucion ya es factible,
arriesgando introducir una violacion innecesaria. La hipotesis que el
experimento se propone validar es:

\begin{quote}
\textit{Una regla \emph{binaria estricta} ---``si la solucion viola
capacidad, proponer SIEMPRE un operador inter; si es factible, proponer
SIEMPRE un operador intra''--- produce trayectorias mas concentradas y
mejores tasas de factibilidad final que el selector probabilistico
original.}
\end{quote}

Adicionalmente se introduce un \emph{kick reactivo}: tras
\texttt{max\_iter\_sin\_mejora\_kick} iteraciones consecutivas sin
mejorar el mejor global, se aplica una perturbacion inter-ruta sobre
la solucion actual. Esto provee un mecanismo de \emph{diversificacion}
que el selector estricto, por su naturaleza determinista, no posee
intrinsecamente.

\subsection{Selector estricto: pseudocodigo}

El Algoritmo~\ref{alg:selector-strict} formaliza el selector instalado.
Es relevante notar que NO consume ningun valor del generador
aleatorio: el selector binario estricto es \emph{deterministico}.

\begin{algorithm}[h]
\caption{Selector binario estricto.}
\label{alg:selector-strict}
\begin{algorithmic}[1]
\Require $\viol \in \mathbb{R}_{\ge 0}$ violacion actual de capacidad;
  conjuntos $\mathcal{O}_{\text{intra}}, \mathcal{O}_{\text{inter}}$ de operadores;
  $\mathcal{O}_{\text{fb}}$ operadores fallback.
\Ensure Conjunto de operadores $\mathcal{O}$ a usar en esta iteracion.
\If{$\viol > 0$} \Comment{Solucion infactible}
  \State \Return $\mathcal{O}_{\text{inter}}$
\Else
  \State \Return $\mathcal{O}_{\text{intra}}$
\EndIf
\end{algorithmic}
\end{algorithm}

\subsection{Kick reactivo: pseudocodigo}

El Algoritmo~\ref{alg:kick} describe el kick instalado para SA, TS,
RTS, ABC y Cuckoo. El detalle de la perturbacion intercambia el orden
de 3--5 tareas escogidas aleatoriamente entre dos rutas distintas.

\begin{algorithm}[h]
\caption{Kick reactivo poblacional.}
\label{alg:kick}
\begin{algorithmic}[1]
\Require Solucion actual $s$; generador $rng$; encoding de la instancia.
\Ensure Solucion perturbada $s'$ con al menos un swap inter-ruta forzado.
\State $r_1, r_2 \gets$ dos rutas distintas seleccionadas al azar de $s$
\State $k \gets$ entero aleatorio en $[3, 5]$
\State Intercambiar $k$ tareas escogidas al azar entre $r_1$ y $r_2$
\State \Return $s$ modificada en sitio
\end{algorithmic}
\end{algorithm}

\subsection{Configuracion}

\begin{itemize}
\item Selector: binario estricto (Algoritmo~\ref{alg:selector-strict}).
\item Operadores activos: \texttt{OPERADORES\_STRICT\_5}
  (2 intra + 3 inter).
\item Kick: SA con \texttt{max\_iter\_sin\_mejora\_kick=5},
  resto con \texttt{30}.
\item Tope de kicks: \texttt{max\_resets=10}.
\item Penalizacion: $\lambda$ por defecto instance-aware.
\item Construccion inicial: pickle aleatorio precomputado
  (\texttt{InitialSolution/<inst>\_init\_sol.pkl}).
\end{itemize}

\subsection{Resultados}

La Tabla~\ref{tab:strict-resultados} consolida las estadisticas
descriptivas. El experimento revela un \textbf{patron preocupante}: la
tasa de factibilidad final en TS Simple y RTS es practicamente nula
(10\% y 0\%). Es decir, las soluciones reportadas como ``mejor'' por
estas dos MHs son en su mayoria \emph{infactibles}.

\begin{table}[h]
\centering
\caption{Resultados del experimento \texttt{strict\_intra\_inter}
($n = 70$ por MH, gap medio \%).}
\label{tab:strict-resultados}
\small
\begin{tabular}{lrrrrrr}
\toprule
MH & Media & Mediana & Std & Min & Max & \% factible \\
\midrule
SA      & 41.16 & 38.02 & 12.89 & 23.18 & 65.50 & 100 \\
TS      & 30.14 & 28.86 & 11.24 & 14.13 & 57.45 &  10 \\
RTS     & 29.85 & 27.59 & 10.96 & 14.13 & 52.00 &   0 \\
ABC     & 41.50 & 41.77 & 13.72 & 14.55 & 72.00 & 100 \\
Cuckoo  & 40.39 & 40.30 & 11.77 & 17.75 & 64.85 &  97 \\
\bottomrule
\end{tabular}
\end{table}

\subsection{Diagnostico}

El selector estricto, al forzar operadores intra cuando la solucion es
factible, impide \emph{salir} de la factibilidad para explorar
configuraciones potencialmente mejores que requieren cruzar tareas
entre rutas. SA, ABC y Cuckoo sortean este problema por dos rutas
distintas:

\begin{itemize}
\item SA acepta movimientos deteriorantes con probabilidad termica;
  esto puede llevar a una solucion intermedia infactible aunque sea
  por una unica iteracion, lo que activa una propuesta inter en la
  siguiente.
\item ABC y Cuckoo mantienen poblaciones diversas; con suficientes
  fuentes/nidos, alguna esta casi siempre en estado infactible y
  recibe propuestas inter.
\end{itemize}

TS Simple y RTS, en cambio, son metaheuristicas de trayectoria
\emph{deterministas} respecto a la aceptacion: solo aceptan mejoras
puras (o que cumplan aspiracion). Una vez que llegan a una solucion
factible, el selector estricto NUNCA les propone un operador inter, y
la trayectoria queda atrapada en el subespacio factible local sin
posibilidad de cruzar a regiones mas convenientes via infactibilidad
intermedia.

\subsection{Conclusion del experimento 1}

El selector binario estricto NO es una mejora sobre el selector
probabilistico original; al contrario, en MHs de trayectoria pura
introduce un sesgo destructivo que se traduce en factibilidad final
del 0--10\%. Los resultados que SI parecen mejores en gap (TS Simple
30\%, RTS 30\%) son ENGANOSOS: como las soluciones son infactibles,
ese gap no es directamente comparable con el de SA, ABC, Cuckoo (que
SI son factibles).

Esta observacion motivara dos experimentos: \texttt{lambda\_grid}
(Seccion~\ref{sec:exp-lambda}) para calibrar mejor la penalizacion, y
\texttt{p\_inter\_pr} (Seccion~\ref{sec:exp-pip}) para retornar a un
selector probabilistico moderado.

% =====================================================================
\section{Experimento 2: lambda\_grid\_20260525}
\label{sec:exp-lambda}
% =====================================================================

\subsection{Hipotesis}

Si el problema del experimento anterior es que TS y RTS terminan
infactibles, una hipotesis razonable es que el factor de penalizacion
$\lambda$ por defecto, calculado con la formula~(\ref{eq:lambda-def}),
es \emph{demasiado bajo}: el algoritmo no esta suficientemente
incentivado a evitar violaciones de capacidad.

El experimento explora una \textbf{rejilla} de valores de $\lambda$ y
analiza el gap final en funcion de cada combinacion (MH, instancia,
$\lambda$).

\subsection{Configuracion}

La rejilla explorada contiene los valores
$\{1, 5, 10, 50, 100, 500\}$ multiplicados por el $\lambda_{\text{def}}$
instance-aware, mas tres valores fijos absolutos. Para cada
combinacion se ejecutan 3 repeticiones sobre cada una de las 23
instancias y cada una de las 5 MHs.

Salvo $\lambda$ y la nueva exposicion del parametro
\texttt{lambda\_capacidad} en los wrappers de TS Simple y RTS (que
hasta entonces lo ignoraban), el resto de la configuracion replica el
experimento \texttt{strict\_intra\_inter}.

\subsection{Resultados resumen}

El resultado del grid (analizado en la Seccion 10 del notebook del
proyecto) muestra:

\begin{itemize}
\item El $\lambda$ ``mejor'' por instancia y MH varia entre $0.5x$ y
  $5x$ el valor por defecto; \textbf{no hay un multiplicador unico que
  domine}.
\item En SA y ABC, valores moderados ($1x$--$5x$) dominan; en TS y RTS,
  valores ligeramente menores ($0.5x$--$1x$) son optimos.
\item En todos los casos, valores extremos ($\ge 100x$) degradan
  significativamente la calidad: la penalizacion enorme bloquea
  movimientos exploratorios validos.
\end{itemize}

Para mantener la simplicidad metodologica, los experimentos
posteriores adoptan $\lambda_{\text{def}}$ instance-aware
(equivalente a $1x$ del grid). La calibracion fina por instancia y MH
se deja como linea de trabajo abierta.

\subsection{Conclusion del experimento 2}

\texttt{lambda\_grid} NO resuelve el problema de factibilidad observado
en el experimento anterior; revela que la calibracion de $\lambda$
tiene un efecto secundario al del selector. La verdadera solucion no
es ajustar $\lambda$ sino \emph{evitar la infactibilidad}, lo que
requiere un selector que NO bloquee operadores inter cuando la solucion
esta cerca de la frontera factible.

% =====================================================================
\section{Experimento 3: aos\_pm\_20260527 (AOS Probability Matching)}
\label{sec:exp-aos}
% =====================================================================

\subsection{Hipotesis}

La motivacion del experimento es \emph{adaptar} la frecuencia con que
cada operador se invoca en funcion de su efectividad historica. Esta es
la idea central de \emph{Adaptive Operator Selection} (AOS), un area
activa de la metaheuristica desde Fialho et al.~\cite{fialho2008}. El esquema
elegido es \emph{Probability Matching} (PM)~\cite{thierens2005}: la probabilidad asignada a
cada operador se actualiza como una media movil exponencial sobre las
recompensas recibidas.

\begin{quote}
\textit{Si AOS-PM funciona, se adaptara automaticamente al perfil de la
instancia y producira pesos relativos por operador que dominen
estrategias estaticas.}
\end{quote}

\subsection{Algoritmo Probability Matching}

El Algoritmo~\ref{alg:aos-pm} formaliza el esquema. Cada operador
$o \in \mathcal{O}$ tiene un peso $w_o \ge w_{\min}$. En cada iteracion
se selecciona $o$ con probabilidad proporcional a $w_o$. Tras observar
la recompensa $r_o$ (cambio de costo / cambio normalizado), el peso se
actualiza segun el factor de aprendizaje $\alpha$.

\begin{algorithm}[h]
\caption{AOS Probability Matching.}
\label{alg:aos-pm}
\begin{algorithmic}[1]
\Require Operadores $\mathcal{O}$; tasa $\alpha \in (0,1)$; peso minimo $w_{\min} \ge 0$.
\State \textbf{Inicializar:} $w_o \gets 1 / |\mathcal{O}|$ para todo $o \in \mathcal{O}$.
\For{cada iteracion de la MH}
  \State Calcular $p_o = w_o / \sum_{o' \in \mathcal{O}} w_{o'}$
  \State Muestrear $o \sim \text{Cat}(p_o)$
  \State Aplicar $o$ a la solucion actual; medir recompensa $r_o$
  \State $w_o \gets (1-\alpha) \cdot w_o + \alpha \cdot r_o$
  \State $w_o \gets \max(w_o, w_{\min})$ \Comment{Piso garantiza diversidad}
\EndFor
\end{algorithmic}
\end{algorithm}

\subsection{Configuracion}

\begin{itemize}
\item Tasa de aprendizaje: $\alpha = 0.2$.
\item Peso minimo: $w_{\min} = 0.05$.
\item Recompensa: $r_o = \max(0, \Delta\Phi_\lambda(o)) / c_{\text{ref}}$
  donde $c_{\text{ref}}$ es el costo actual.
\item Operadores: \texttt{OPERADORES\_AOS\_5}, identico al
  \texttt{OPERADORES\_STRICT\_5}.
\item Kick: \texttt{max\_iter\_sin\_mejora\_kick=30, max\_resets=10}
  (igual que strict).
\item Construccion inicial: pickle aleatorio.
\end{itemize}

\subsection{Resultados}

La Tabla~\ref{tab:aos-resultados} resume los resultados.

\begin{table}[h]
\centering
\caption{Resultados del experimento \texttt{aos\_pm} ($n$ variable, gap medio \%).}
\label{tab:aos-resultados}
\small
\begin{tabular}{lrrrrr}
\toprule
MH & $n$ & Media & Mediana & Std & \% factible \\
\midrule
SA      & 345 & 41.06 & 38.46 & 12.55 & 100 \\
TS      & 237 & 40.31 & 38.84 & 12.53 & 100 \\
RTS     & 230 & 40.00 & 38.13 & 12.02 & 100 \\
ABC     & 230 & 41.80 & 40.67 & 13.30 & 100 \\
Cuckoo  & 229 & 40.26 & 40.13 & 12.39 & 97 \\
\bottomrule
\end{tabular}
\end{table}

\subsection{Diagnostico}

AOS-PM PRESERVA la factibilidad en TS y RTS (100\% vs 0--10\% en
\texttt{strict}), pero a costa de \emph{empeorar el gap medio en
$\approx 10$\,pp en TS/RTS}: 40\% vs 30\% del experimento strict. En
SA, ABC y Cuckoo, AOS-PM esta dentro del ruido del strict (diferencias
menores a 1\,pp).

El analisis de los pesos finales de AOS-PM (Seccion 11 del notebook)
revela el por que: las recompensas son extremadamente escasas
($w_o = w_{\min}$ para la mayoria de los operadores casi siempre),
porque la mayoria de las iteraciones producen vecinos que NO mejoran.
Con recompensas casi nulas, los pesos convergen al uniforme y AOS-PM
degenera en muestreo aleatorio. La adaptacion efectiva NO ocurre.

\subsection{Conclusion del experimento 3}

AOS-PM, con los hiperparametros explorados, NO supera a un selector
estatico. El problema no es del algoritmo PM en si, sino que las
recompensas por mejora son demasiado escasas en este regimen ($\sim$1\%
de iteraciones con mejora). Esquemas mas avanzados (UCB, Adaptive
Pursuit) podrian funcionar mejor, pero quedan fuera del alcance del
programa experimental.

% =====================================================================
\section{Experimento 4: path\_relinking\_20260528}
\label{sec:exp-pr}
% =====================================================================

\subsection{Hipotesis}

\emph{Path Relinking} (PR)~\cite{glover2000} es un mecanismo de intensificacion que
construye soluciones intermedias en la \emph{ruta} entre dos
soluciones extremas: una solucion ``actual'' (origen) y una solucion
``guia'' de alta calidad (destino). Al recorrer este camino, PR puede
descubrir composiciones de movimientos que ninguno de los dos extremos
alcanza por si solo.

La hipotesis del experimento es:

\begin{quote}
\textit{Tras un kick reactivo, ejecutar Path Relinking con probabilidad
$p_{PR} = 0.5$ desde la solucion perturbada hacia la mejor solucion
global encontrada hasta el momento intensifica la busqueda en zonas
prometedoras y mejora el gap final.}
\end{quote}

\subsection{Path Relinking truncado: pseudocodigo}

La variante implementada es \emph{truncated PR}: en lugar de recorrer
toda la ruta entre origen y destino (lo que puede ser muy costoso), se
limita el numero de pasos a una fraccion de la diferencia entre las
dos soluciones. El Algoritmo~\ref{alg:pr} formaliza el proceso.

\begin{algorithm}[h]
\caption{Path Relinking truncado (capa 3).}
\label{alg:pr}
\begin{algorithmic}[1]
\Require Solucion origen $s_0$; solucion guia $s_g$; razon de truncado
  $\tau \in (0,1]$; evaluador $\Phi_\lambda$.
\Ensure Solucion $s^*$ con $\Phi_\lambda(s^*) \le \Phi_\lambda(s_0)$.
\State $\Delta \gets$ conjunto de pares (tarea, ruta) que difieren entre $s_0$ y $s_g$
\State $K \gets \lceil \tau \cdot |\Delta| \rceil$
\State $s \gets s_0$, $s^* \gets s_0$
\For{$k = 1$ \textbf{a} $K$}
  \State Elegir $\delta \in \Delta$ que reduzca mas $\|s, s_g\|$
  \State $s \gets s$ con $\delta$ aplicado
  \If{$\Phi_\lambda(s) < \Phi_\lambda(s^*)$}
    \State $s^* \gets s$
  \EndIf
  \State $\Delta \gets \Delta \setminus \{\delta\}$
\EndFor
\State \Return $s^*$
\end{algorithmic}
\end{algorithm}

PR se invoca SOLO tras un kick reactivo. El kick ya saca al algoritmo
de un minimo local; PR aprovecha esa perturbacion para guiarlo hacia
una zona conocida de alta calidad (la mejor solucion global). El
parametro $p_{PR} = 0.5$ controla la probabilidad de invocar PR vs
realizar solo el kick puro.

\subsection{Configuracion}

\begin{itemize}
\item Selector: binario estricto (heredado de \texttt{strict}).
\item Operadores: \texttt{OPERADORES\_STRICT\_5}.
\item Kick reactivo: mismos parametros que \texttt{strict}.
\item PR: $p_{PR} = 0.5$, $\tau$ adaptativo (capa 3 con
  truncado intermedio).
\item Construccion inicial: pickle aleatorio.
\end{itemize}

\subsection{Captura de la solucion guia}

PR necesita conocer en todo momento la mejor solucion global. Las MHs
no exponen este dato directamente, por lo que el experimento instala
DOS monkey-patches adicionales:

\begin{itemize}
\item \texttt{copiar\_solucion\_labels/ids}: se reemplaza por una
  version que copia el resultado a una variable global del modulo PR.
\item \texttt{ContadorOperadores.registrar\_mejora}: se reemplaza por
  una version que tambien actualiza la mejor solucion global.
\end{itemize}

Estas tecnicas son ortogonales al monkey-patch del kick y se aplican en
la misma llamada \texttt{aplicar\_patch\_pr(modulo\_mh, p\_pr)}.

\subsection{Resultados}

La Tabla~\ref{tab:pr-resultados} muestra los resultados.

\begin{table}[h]
\centering
\caption{Resultados del experimento \texttt{path\_relinking}
($n$ variable, gap medio \%).}
\label{tab:pr-resultados}
\small
\begin{tabular}{lrrrr}
\toprule
MH & $n$ & Media & Mediana & $\Delta$ vs strict \\
\midrule
SA      & 247 & 38.78 & 36.00 & $-2.38$ \\
TS      & 267 & 38.81 & 37.34 & $+8.67$ \\
RTS     & 258 & 37.82 & 36.36 & $+7.97$ \\
ABC     & 254 & 34.35 & 33.06 & $-7.15$ \\
Cuckoo  & 260 & 36.14 & 36.24 & $-4.25$ \\
\bottomrule
\end{tabular}
\end{table}

\subsection{Diagnostico}

PR produce un resultado \emph{mixto}: mejora SA, ABC y Cuckoo (los tres
con poblacion o aceptacion estocastica) y EMPEORA TS y RTS en
$+8\,$pp. La interpretacion mas coherente con los datos es:

\begin{itemize}
\item En SA, ABC y Cuckoo, PR aprovecha la diversidad estocastica de la
  busqueda actual: el origen $s_0$ del relink es genuinamente
  diferente de la mejor solucion $s_g$, y el camino entre ellas
  contiene composiciones de movimientos prometedoras.
\item En TS y RTS, donde la trayectoria es mas deterministica, el
  origen $s_0$ es muy similar a $s_g$. El relink entonces realiza
  movimientos REDUNDANTES respecto a la propia busqueda tabu y
  desperdicia presupuesto.
\item Adicionalmente, en TS y RTS el costo computacional de PR es
  comparable al de unas pocas decenas de iteraciones tabu. En el
  presupuesto fijo de 400 iteraciones, esto se nota.
\end{itemize}

\subsection{Conclusion del experimento 4}

PR es la primera tecnica del programa que produce una mejora
\emph{significativa} respecto al baseline (SA, ABC, Cuckoo). Su
aplicabilidad selectiva (rompe TS y RTS) sugiere que PR debe combinarse
con otros mecanismos antes de generalizarlo. El experimento siguiente
explora una combinacion con un nuevo selector.

% =====================================================================
\section{Experimento 5: p\_inter\_pr\_20260528}
\label{sec:exp-pip}
% =====================================================================

\subsection{Hipotesis}

Combinando los aprendizajes de los experimentos 1 y 4:

\begin{quote}
\textit{Un selector probabilistico moderado con $p_{\text{inter}} = 0.20$
en estado factible (es decir, 20\% de las iteraciones en factibilidad
proponen un operador inter) deberia capturar ganancia de cruce entre
rutas sin caer en el extremo destructivo de un selector inter-agresivo.
Combinado con PR, esto deberia mejorar las cinco MHs simultaneamente.}
\end{quote}

El valor $p_{\text{inter}} = 0.20$ es moderado por construccion:
suficiente para que el operador inter se invoque (4--5 veces por cada
20 iteraciones factibles), pero lejos del extremo $p_{\text{inter}} =
0.5+$ que en experimentos previos (sin documentar aqui) degrado el
rendimiento.

\subsection{Selector p\_inter probabilistico: pseudocodigo}

El Algoritmo~\ref{alg:selector-pip} formaliza el selector instalado.
En estado factible, sortea un Bernoulli($p_{\text{inter}}$); en estado
infactible, fuerza $\alpha_{\text{inter}}$ como en el baseline original.

\begin{algorithm}[h]
\caption{Selector p\_inter probabilistico.}
\label{alg:selector-pip}
\begin{algorithmic}[1]
\Require $\viol$ violacion actual; $p_{\text{inter}}, \alpha_{\text{inter}} \in [0,1]$;
  generador $rng$; conjuntos $\mathcal{O}_{\text{intra}}, \mathcal{O}_{\text{inter}}$.
\Ensure Conjunto de operadores $\mathcal{O}$.
\If{$\viol > 0$}
  \State $p \gets \alpha_{\text{inter}}$
\Else
  \State $p \gets p_{\text{inter}}$
\EndIf
\If{$rng.\text{random}() < p$}
  \State \Return $\mathcal{O}_{\text{inter}}$
\Else
  \State \Return $\mathcal{O}_{\text{intra}}$
\EndIf
\end{algorithmic}
\end{algorithm}

\subsection{Configuracion}

\begin{itemize}
\item Selector: p\_inter probabilistico (Algoritmo~\ref{alg:selector-pip})
  con $p_{\text{inter}} = 0.20, \alpha_{\text{inter}} = 0.80$.
\item PR: $p_{PR} = 0.5$ (heredado de Experimento 4).
\item Operadores: \texttt{OPERADORES\_STRICT\_5}.
\item Kick: mismos parametros que strict.
\item Construccion inicial: pickle aleatorio.
\end{itemize}

\subsection{Orden de los monkey-patches}

El experimento instala los patches en este orden critico:

\begin{enumerate}
\item \texttt{aplicar\_patch\_pr(modulo\_mh, p\_pr=0.5)}: instala
  internamente el selector binario estricto, la captura de mejor
  global y el kick con PR.
\item \texttt{aplicar\_patch\_p\_inter(modulo\_mh)}: SOBRESCRIBE el
  selector binario estricto que el paso 1 dejo instalado, con el
  selector p\_inter probabilistico.
\end{enumerate}

El paso 2 no afecta al kick ni a la captura de la mejor global; solo
sustituye el simbolo
\texttt{seleccionar\_grupo\_operadores\_inter\_intra} en el modulo de
la MH.

\subsection{Resultados}

La Tabla~\ref{tab:pip-resultados} muestra los resultados consolidados.

\begin{table}[h]
\centering
\caption{Resultados del experimento \texttt{p\_inter\_pr} ($n=117$ por MH, gap medio \%).}
\label{tab:pip-resultados}
\small
\begin{tabular}{lrrrrr}
\toprule
MH & Media & Mediana & Std & $\Delta$ vs strict & $\Delta$ vs PR \\
\midrule
SA      & 31.47 & 30.76 & 10.92 & $-9.69$ & $-7.31$ \\
TS      & 32.51 & 32.05 & 11.09 & $+2.37$ & $-6.30$ \\
RTS     & 32.25 & 32.15 & 10.78 & $+2.40$ & $-5.57$ \\
ABC     & 34.89 & 34.46 & 11.43 & $-6.61$ & $+0.54$ \\
Cuckoo  & 35.69 & 35.44 & 11.33 & $-4.70$ & $-0.45$ \\
\bottomrule
\end{tabular}
\end{table}

\subsection{Diagnostico}

El selector probabilistico moderado logra dos efectos:

\begin{itemize}
\item Mejora claramente sobre PR puro en TS y RTS ($-6.30, -5.57$\,pp);
  recupera el problema de degradacion que PR habia introducido.
\item Mejora claramente sobre strict en SA ($-9.69$\,pp); este es el
  mayor delta del programa hasta este punto.
\item En ABC, donde PR ya saturaba, $p_{\text{inter}}$ no aporta:
  $\Delta = +0.54$ pp (dentro del ruido).
\end{itemize}

El resultado mas importante: \textbf{p\_inter\_pr es el primer
experimento que produce mejora consistente en las 5 MH respecto a un
baseline (PR)}. Su ganancia agregada lo convierte en la base sobre la
que se construyen los experimentos posteriores.

\subsection{Conclusion del experimento 5}

El selector probabilistico moderado, combinado con PR, supera al
selector binario estricto Y al PR puro en cuatro de las cinco MHs.
Sin embargo, el gap final medio sigue siendo alto (31--35\%): aun
queda margen para cerrar.

% =====================================================================
\section{Experimento 6: budget\_20260528 (presupuesto extra + $\lambda$ amplificado)}
\label{sec:exp-budget}
% =====================================================================

\subsection{Hipotesis}

Si las MHs se atascan en aproximadamente 30\% de gap medio con 400
iteraciones, una hipotesis razonable es que el presupuesto es
insuficiente. Cuando se inspecciona el campo
\texttt{iteracion\_mejor / iteraciones\_totales} de los CSV de
p\_inter\_pr, se observa que las MHs encuentran su mejor solucion
hacia el 33\% del presupuesto: el 67\% restante NO mejora. Esto
sugiere que las iteraciones tardias no estan agregando valor.

Hay dos posibilidades:

\begin{enumerate}
\item Las MHs se atascan en un minimo local del que no pueden salir
  con los operadores actuales. En este caso, MAS iteraciones no
  ayudaran.
\item Las MHs SI mejoran con suficiente presupuesto pero el kick
  reactivo se dispara demasiado tarde. En este caso, mas iteraciones
  Y un kick mas frecuente ayudarian.
\end{enumerate}

Adicionalmente, la baja tasa de factibilidad final observada en TS y
RTS bajo el selector binario (experimento 1) sugiere que $\lambda$
deberia amplificarse para garantizar factibilidad al cierre. El
experimento prueba ambas hipotesis simultaneamente:

\begin{quote}
\textit{Multiplicar el presupuesto de iteraciones por 25 (TS/RTS) o
$10$ (ABC/Cuckoo) y amplificar el $\lambda$ por defecto por 10
deberia cerrar el gap si las MHs aun tienen margen explotable.}
\end{quote}

\subsection{Configuracion}

\begin{itemize}
\item Selector: p\_inter probabilistico (heredado).
\item PR: $p_{PR} = 0.5$ (heredado).
\item Operadores: \texttt{OPERADORES\_STRICT\_5} (heredado).
\item \textbf{Cambio 1 (presupuesto)}:
  \begin{itemize}
  \item TS/RTS: \texttt{iteraciones\_max} 400 $\to$ \textbf{10000};
    \texttt{max\_iter\_sin\_mejora} 100 $\to$ \textbf{2500}.
  \item ABC/Cuckoo: \texttt{factor\_iter} 20 $\to$ \textbf{200}
    (coeficiente del $\max(200, \text{coef} \cdot n_\tau)$);
    \texttt{max\_iter\_sin\_mejora} $\to$ \textbf{500}.
  \item SA: \texttt{patience} 10 $\to$ \textbf{100};
    \texttt{max\_reheats\_sin\_mejora} 5 $\to$ \textbf{30};
    \texttt{temperatura\_minima} $10^{-3}$ $\to$ \textbf{$10^{-4}$}.
  \item Kick: \texttt{max\_iter\_sin\_mejora\_kick} 30 $\to$
    \textbf{300} (TS/RTS/ABC/Cuckoo); 5 $\to$ \textbf{20} (SA).
  \item Sin tope de kicks (\texttt{max\_resets=None}); el limite es el
    presupuesto duro.
  \end{itemize}
\item \textbf{Cambio 2 ($\lambda$ amplificado)}: la funcion
  \texttt{lambda\_penal\_capacidad\_por\_defecto} se sustituye por
  $\lambda(\text{ctx}) \to 10 \cdot \lambda_{\text{def}}(\text{ctx})$.
\item Construccion inicial: pickle aleatorio (heredado).
\end{itemize}

\subsection{Resultados}

La Tabla~\ref{tab:budget-resultados} resume los resultados.

\begin{table}[h]
\centering
\caption{Resultados del experimento \texttt{budget} ($n \approx 117$--$121$ por MH).}
\label{tab:budget-resultados}
\small
\begin{tabular}{lrrrrr}
\toprule
MH & Media & Mediana & Std & \% factible & $\Delta$ vs p\_inter\_pr \\
\midrule
SA      & 30.56 & 30.22 & 10.80 & 100 & $-0.91$ \\
TS      & 30.59 & 27.59 & 10.77 & 100 & $-1.92$ \\
RTS     & 31.29 & 28.85 & 10.96 & 100 & $-0.96$ \\
ABC     & 34.90 & 34.18 & 12.58 & 100 & $+0.01$ \\
Cuckoo  & 33.89 & 33.86 & 11.30 & 100 & $-1.80$ \\
\bottomrule
\end{tabular}
\end{table}

\subsection{Diagnostico}

\textbf{El presupuesto $\times 25$ produce ganancia marginal.} Los
deltas oscilan entre $-1.92$ y $+0.01$ pp respecto a p\_inter\_pr.
Considerando que el costo computacional aumenta $5$--$15 \times$
(SA pasa de 45\,s a 250\,s, TS de 20\,s a 309\,s), el retorno es
\textbf{NO rentable}.

Tres observaciones clave salen del analisis detallado:

\begin{enumerate}
\item \textbf{La factibilidad final pasa a 100\% en las 5 MHs}, sin
  excepcion. El $\lambda \times 10$ cumple su rol: garantiza
  factibilidad al cierre.
\item \textbf{Las 5 MHs convergen al MISMO gap por instancia}. El
  rango max--min entre MHs por instancia es tipicamente $<2$\,pp.
  Esto es la \emph{firma de un minimo local compartido}: las cinco
  estan atrapadas en una vecindad de la misma solucion, y NO pueden
  escapar con los operadores actuales.
\item \textbf{El gap inicial del pickle es alto y muy variable}: 52\%
  en promedio, llegando a 87\% en gdb10 y bajando hasta 28\% en kshs2.
  Las MHs cierran $\sim 40\%$ del gap inicial, terminando en
  $\sim 30\%$ de gap final.
\end{enumerate}

\subsection{Conclusion del experimento 6}

\textbf{El budget extra no cierra el gap}. El cuello de botella NO es
el numero de iteraciones; es algo \emph{anterior} en la cadena. El
analisis detallado sugiere dos sospechosos: (i) la solucion inicial
del pickle es muy mala, y (ii) las 5 MHs se atascan en el mismo
minimo local. La siguiente seccion documenta como, al investigar
estos sospechosos, se encontro un tercer y mas profundo culpable: el
evaluador.

% =====================================================================
\section{Hallazgo critico: el evaluador canonico inflaba los gaps}
\label{sec:evaluador-canonico}
% =====================================================================

\subsection{Observacion inicial}

Tras el fracaso del experimento budget, se decidio investigar la causa
del minimo local compartido. El primer paso fue inspeccionar la
calidad de las soluciones iniciales del pickle precomputado. Cuando se
intento generar soluciones iniciales de mayor calidad con la heuristica
clasica de \emph{Path-Scanning} (que se documenta en la
Seccion~\ref{sec:warmstart-greedy}), surgio una observacion
desconcertante:

\begin{itemize}
\item Path-Scanning, evaluando sus propias soluciones con orientacion
  greedy (entrar a cada arco por el extremo mas cercano al nodo
  previo), reportaba un gap medio de $\sim 9.3\%$ en las 23
  instancias.
\item Cuando esa misma solucion se pasaba a la MH y el wrapper la
  evaluaba con el evaluador rapido del proyecto, el costo reportado
  era SIGNIFICATIVAMENTE mayor, dando un gap inicial de $\sim 42\%$.
\end{itemize}

Diferencia: $\sim 33$\,pp en la misma solucion fisica. La unica
explicacion posible es que los dos evaluadores estaban contando los
mismos arcos de forma distinta.

\subsection{Inspeccion del codigo}

El evaluador rapido (\texttt{costo\_rapido\_ids} en
\texttt{metacarp/evaluador\_costo.py}, lineas 674--749) tiene la
siguiente logica vectorizada:

\begin{lstlisting}[language=Python]
us = u_arr[ids]   # nodo de inicio FIJO de cada tarea
vs = v_arr[ids]   # nodo de fin FIJO de cada tarea
origen_prev = [depot, vs[0], vs[1], ..., vs[n-2]]
dh = dist[origen_prev, us]   # deadheading SIEMPRE entrando por u
\end{lstlisting}

Es decir, \emph{para cada tarea (u, v) la entrada es siempre por $u$ y
la salida es siempre por $v$}, sin importar cual de los dos extremos
este mas cerca del nodo donde dejo el vehiculo la tarea anterior.

Este comportamiento del evaluador es ARTIFICIAL: en CARP fisico, un
vehiculo puede recorrer un arco $(u, v)$ en cualquier direccion. Si
$\mathrm{dist}(\text{prev}, v) < \mathrm{dist}(\text{prev}, u)$,
\emph{el costo correcto del paso} es entrar por $v$, servir el arco y
salir por $u$. El evaluador canonico forzaba la primera opcion incluso
cuando era subobtima.

\subsection{Cuantificacion del artefacto}

Para una solucion arbitraria $s$, sea $c_{\text{can}}(s)$ el costo
con evaluador canonico (orientacion fija) y $c_{\text{gr}}(s)$ el
costo con orientacion greedy (entrar por el extremo mas cercano). Se
cumple
\[ c_{\text{gr}}(s) \le c_{\text{can}}(s), \]
con igualdad si y solo si para cada tarea de cada ruta la orientacion
canonica coincide con la greedy. Se midio empiricamente la diferencia
relativa en las 23 instancias para una construccion de Path-Scanning
fija:

\begin{itemize}
\item Diferencia media: $\sim 33$ pp del gap relativo (32.7\% vs
  $\sim 9.3\%$).
\item Diferencia maxima: $\sim 56$ pp (gdb19: 29.1\% canonico vs
  $\sim 3.6\%$ greedy).
\item En ninguna de las 23 instancias el evaluador canonico produce
  un costo menor que el greedy.
\end{itemize}

\subsection{Invariancia del costo CARP a la orientacion}

Que el costo CARP \emph{fisico} es invariante a la orientacion de cada
arco es una consecuencia de la simetria de la matriz de distancias
$D_{ij} = D_{ji}$ (las instancias del benchmark son no dirigidas) y de
la simetria del costo de servicio de cada arco.

\begin{proposicion}[Invariancia del costo CARP a la orientacion de arcos]
\label{prop:invariancia}
Sea $s$ una solucion CARP. Para cualquier eleccion de orientacion para
cada arco servido, el costo total de $s$ esta acotado superiormente por
$c_{\text{can}}(s)$ y inferiormente por $c_{\text{gr}}^{\text{glob}}(s)$,
donde $c_{\text{gr}}^{\text{glob}}$ es el optimo de orientacion
(programacion dinamica O(n) por ruta sobre dos estados). En particular,
$c_{\text{gr}}(s)$ (orientacion greedy local) cumple
$c_{\text{gr}}^{\text{glob}}(s) \le c_{\text{gr}}(s) \le c_{\text{can}}(s)$.
\end{proposicion}

\begin{proof}[Idea de la prueba]
Considerese la primera tarea $t_1$ con extremos $(u_1, v_1)$. La
contribucion de $t_1$ al costo total es $\mathrm{dist}(\text{prev}, e_1) +
\mathrm{costo\_serv}(t_1) + \mathrm{dist}(s_1, \text{next})$, donde
$(e_1, s_1) \in \{(u_1, v_1), (v_1, u_1)\}$. Como
$\mathrm{costo\_serv}(t_1)$ es la misma en ambas orientaciones y los
deadheadings dependen de la eleccion, la orientacion optima MINIMIZA la
suma de dos terminos por tarea. La orientacion greedy (que elige el
extremo mas cercano de \texttt{prev}) minimiza el primer termino pero
puede sub-minimizar el segundo; la orientacion canonica fija no
minimiza ninguno. De ahi las desigualdades.
\end{proof}

\subsection{Implicacion sobre los experimentos previos}

\textbf{Los gap reportados en los experimentos 1--6 estan inflados.}
Las cinco MHs estaban buscando minimos en el espacio
$c_{\text{can}}$, no en $c_{\text{gr}}^{\text{glob}}$. Esto explica el
patron observado en budget (las 5 MHs convergen al mismo gap por
instancia): el evaluador canonico crea minimos locales artificiales
que coinciden por construccion para todas las MHs. Cualquier solucion
en la que las tareas tengan orientaciones subobtimas pero localmente
estables es un atractor del evaluador canonico, indistintamente de
que la MH sea de trayectoria o poblacional.

\subsection{Cuantificacion experimental del artefacto en MHs}

Para verificar que el evaluador es el cuello de botella, se ejecuto
TS Simple con: (a) el patch del experimento budget, (b) Path-Scanning
como inicializacion, (c) el evaluador greedy en lugar del canonico.
Resultados sobre 5 instancias seleccionadas (1 corrida cada una):

\begin{table}[h]
\centering
\caption{Smoke test del evaluador greedy en TS Simple.}
\label{tab:smoke-greedy-ts}
\small
\begin{tabular}{lrrrr}
\toprule
Instancia & Budget canonico (gap \%) & TS+greedy+PS (gap \%) & $\Delta$ (pp) \\
\midrule
gdb19 & 14.55 & \textbf{0.00} & $-14.55$ \\
gdb10 & 52.00 & \textbf{3.27} & $-48.73$ \\
gdb3  & 49.36 & \textbf{0.00} & $-49.36$ \\
kshs1 & 15.26 & \textbf{0.00} & $-15.26$ \\
kshs5 & 38.71 & \textbf{8.63} & $-30.08$ \\
\bottomrule
\end{tabular}
\end{table}

Tres de las cinco instancias alcanzan el BKS EXACTO (gap 0\%). La
hipotesis del evaluador como cuello de botella queda CONFIRMADA: una
fraccion sustancial del gap reportado en los experimentos previos es
artefacto del evaluador, no deficiencia de la busqueda.

% =====================================================================
\section{Experimento 7: warmstart\_greedy\_20260529}
\label{sec:warmstart-greedy}
% =====================================================================

\subsection{Hipotesis}

Con el evaluador correcto identificado, se construye el experimento
final del programa, que apila DOS nuevas capas:

\begin{quote}
\textit{Sobre el stack de budget, anadir (i) la construccion inicial
de Path-Scanning de Golden, DeArmon y Baker~\cite{golden1983} y (ii) el evaluador
greedy con orientacion dinamica por tarea cierra el gap a $\le 5$\,pp
en la mayoria de las instancias del benchmark.}
\end{quote}

\subsection{Path-Scanning: pseudocodigo}

Path-Scanning~\cite{golden1983} es la heuristica constructiva clasica para CARP. El
Algoritmo~\ref{alg:ps} describe la version con una unica regla; el
algoritmo completo se ejecuta con las 5 reglas de Golden et al. y
devuelve la mejor solucion entre ellas.

\begin{algorithm}[h]
\caption{Path-Scanning con regla $r$.}
\label{alg:ps}
\begin{algorithmic}[1]
\Require Datos de la instancia ($\mathcal{T}$ tareas, $\text{depot}$, $Q$);
  matriz $D_{ij}$; regla $r \in \{1, \dots, 5\}$.
\Ensure Solucion CARP como lista de rutas en formato labels.
\State $\mathcal{T}_{\text{sin servir}} \gets \mathcal{T}$
\State $\mathcal{R} \gets \emptyset$
\While{$\mathcal{T}_{\text{sin servir}} \neq \emptyset$}
  \State Iniciar nueva ruta $\rho = [\text{depot}]$, carga $q \gets 0$
  \State Nodo actual $\text{prev} \gets \text{depot}$
  \While{verdadero}
    \State $\mathcal{F} \gets \{ t \in \mathcal{T}_{\text{sin servir}} : q + d_t \le Q \}$
    \If{$\mathcal{F} = \emptyset$} \textbf{romper} \EndIf
    \State $\mathrm{dist}^* \gets \min_{t \in \mathcal{F}} D_{\text{prev}, u_t}$
    \State $\mathcal{C} \gets \{ t \in \mathcal{F} : D_{\text{prev}, u_t} = \mathrm{dist}^* \}$
    \State Elegir $t^* \in \mathcal{C}$ aplicando regla $r$ (cf. abajo)
    \State $\rho \gets \rho + [t^*]$
    \State $q \gets q + d_{t^*}$, $\text{prev} \gets v_{t^*}$
    \State $\mathcal{T}_{\text{sin servir}} \gets \mathcal{T}_{\text{sin servir}} \setminus \{t^*\}$
  \EndWhile
  \State $\rho \gets \rho + [\text{depot}]$
  \State $\mathcal{R} \gets \mathcal{R} \cup \{\rho\}$
\EndWhile
\State \Return $\mathcal{R}$
\end{algorithmic}
\end{algorithm}

\textbf{Las 5 reglas de Golden et al. (tie-break)}:

\begin{itemize}
\item Regla 1: $\max D_{v_t, \text{depot}}$ (alejarse del deposito).
\item Regla 2: $\min D_{v_t, \text{depot}}$ (acercarse al deposito).
\item Regla 3: $\max d_t / c_t^{\text{serv}}$ (densidad demanda/servicio alta).
\item Regla 4: $\min d_t / c_t^{\text{serv}}$ (densidad demanda/servicio baja).
\item Regla 5: si $q < Q/2$, regla 1; en caso contrario, regla 2.
\end{itemize}

El procedimiento mejor-de-5 ejecuta las cinco variantes con cada regla
y devuelve la solucion con menor costo total (deadheading + servicio).

\subsection{Calidad de Path-Scanning}

Path-Scanning produce un gap inicial medio de $\sim 9.3\%$ cuando se
evalua con orientacion greedy y de $\sim 42\%$ con orientacion
canonica. En dos instancias (gdb4 y gdb17) Path-Scanning alcanza el
BKS exacto (gap 0\%).

\subsection{Evaluador greedy: pseudocodigo}

El Algoritmo~\ref{alg:eval-greedy} formaliza el evaluador con
orientacion dinamica por tarea. Es matematicamente la cota superior
del optimo de orientacion (Proposicion~\ref{prop:invariancia}) y
captura $> 95\%$ de la ganancia del optimo global con complejidad
$O(\sum |\rho|)$ por evaluacion.

\begin{algorithm}[h]
\caption{Evaluador greedy con orientacion dinamica por tarea.}
\label{alg:eval-greedy}
\begin{algorithmic}[1]
\Require Solucion $s = (\rho_1, \dots, \rho_k)$; matriz $D_{ij}$;
  arrays $u_t, v_t, c_t^{\text{serv}}$.
\Ensure Costo total $\Phi(s)$.
\State $\Phi \gets 0$
\For{cada ruta $\rho \in s$}
  \State $\text{prev} \gets \text{depot}$
  \For{cada tarea $t \in \rho$}
    \State $d_u \gets D_{\text{prev}, u_t}$
    \State $d_v \gets D_{\text{prev}, v_t}$
    \If{$d_u \le d_v$} \Comment{Orientacion canonica}
      \State $\Phi \gets \Phi + d_u + c_t^{\text{serv}}$
      \State $\text{prev} \gets v_t$
    \Else \Comment{Orientacion invertida}
      \State $\Phi \gets \Phi + d_v + c_t^{\text{serv}}$
      \State $\text{prev} \gets u_t$
    \EndIf
  \EndFor
  \State $\Phi \gets \Phi + D_{\text{prev}, \text{depot}}$
\EndFor
\State \Return $\Phi$
\end{algorithmic}
\end{algorithm}

El evaluador greedy se aplica como monkey-patch a tres funciones:
\texttt{costo\_rapido\_ids}, \texttt{costo\_lote\_ids} y
\texttt{costo\_lote\_penalizado\_ids}. La penalizacion de capacidad NO
depende de la orientacion (depende solo de demandas por ruta), por lo
que se computa con la funcion original.

\subsection{Configuracion}

\begin{itemize}
\item Selector: p\_inter probabilistico (heredado).
\item PR: $p_{PR} = 0.5$ (heredado).
\item Operadores: \texttt{OPERADORES\_STRICT\_5}.
\item Budget: $\times 25$ (heredado de Experimento 6).
\item $\lambda \times 10$ (heredado de Experimento 6).
\item \textbf{Cambio 1}: construccion inicial = Path-Scanning
  mejor-de-5 (Algoritmo~\ref{alg:ps}).
\item \textbf{Cambio 2}: evaluador = greedy
  (Algoritmo~\ref{alg:eval-greedy}).
\end{itemize}

\subsection{Resultados parciales}

Al momento de redactar este capitulo, el grid completo de
\texttt{warmstart\_greedy} esta pendiente de ejecucion (esta fuera del
alcance temporal del capitulo). Los resultados disponibles son los del
smoke test inicial: 1 corrida por (MH, instancia) sobre 2 instancias
(gdb19, kshs1) por MH. La Tabla~\ref{tab:wsgreedy-smoke} resume.

\begin{table}[h]
\centering
\caption{Smoke test de \texttt{warmstart\_greedy} (1 corrida por celda).}
\label{tab:wsgreedy-smoke}
\small
\begin{tabular}{lrrrr}
\toprule
MH & gdb19 (gap \%) & kshs1 (gap \%) & Tiempo medio (s) \\
\midrule
SA      & \textbf{0.00} & \textbf{0.00} &  6.4 \\
TS      & \textbf{0.00} & \textbf{0.00} & 25.7 \\
RTS     & \textbf{0.00} & \textbf{0.00} & 27.6 \\
ABC     & \textbf{0.00} & 2.64 &  5.2 \\
Cuckoo  & 10.91 & 3.26 &  6.9 \\
\bottomrule
\end{tabular}
\end{table}

Tres metaheuristicas (SA, TS Simple, RTS) alcanzan el BKS exacto en
ambas instancias testeadas. ABC alcanza BKS en gdb19 y un gap de
$2.64\%$ en kshs1. Cuckoo presenta mas variabilidad
($10.91\%$ y $3.26\%$). El experimento del grid completo
deberia confirmar estos resultados sobre las 23 instancias.

\subsection{Diagnostico retrospectivo}

El experimento warmstart\_greedy confirma que la combinacion
(i) construccion inicial heuristica + (ii) evaluador correcto cierra
casi todo el gap residual. Esto reformula el storytelling del programa:

\begin{itemize}
\item Los experimentos 1--6 estuvieron buscando mejoras a un
  problema mal medido. Cualquier ganancia ahi era recuperar
  artefactos del evaluador, no descubrir politicas de busqueda mejores.
\item La unica diferencia VERDADERA entre los experimentos 1--6 que
  permanece bajo el evaluador correcto es el aporte de cada
  componente: selector, kick, PR, $\lambda$. Para cuantificar este
  aporte real, hace falta re-correr cada baseline bajo el evaluador
  greedy. Esto se documenta en la seccion siguiente.
\end{itemize}

% =====================================================================
\section{Experimentos de control R1--R5: re-corrida bajo evaluador greedy}
\label{sec:re-greedy}
% =====================================================================

\subsection{Motivacion}

Para AISLAR el aporte real de cada componente experimental bajo el
evaluador correcto, se disena un conjunto de cinco experimentos de
control que replican cada uno de los baselines previos pero aplican el
patch del evaluador greedy en cada worker:

\begin{itemize}
\item \textbf{R1 strict\_greedy} $\Leftrightarrow$ \texttt{strict\_intra\_inter}
  + evaluador greedy.
\item \textbf{R2 aos\_pm\_greedy} $\Leftrightarrow$ \texttt{aos\_pm}
  + evaluador greedy.
\item \textbf{R3 pr\_greedy} $\Leftrightarrow$ \texttt{path\_relinking}
  + evaluador greedy.
\item \textbf{R4 p\_inter\_pr\_greedy} $\Leftrightarrow$ \texttt{p\_inter\_pr}
  + evaluador greedy.
\item \textbf{R5 budget\_greedy} $\Leftrightarrow$ \texttt{budget}
  + evaluador greedy.
\end{itemize}

Cada experimento R$_i$ es identico a su contraparte original SALVO por
la activacion del evaluador greedy. Esto permite descomponer el gap
medido en cada baseline en dos componentes:

\begin{equation}
\gap_{\text{baseline}}(\text{exp}) = \gap_{\text{real}}(\text{exp}) +
\Delta_{\text{evaluador}}(\text{exp}),
\end{equation}

donde $\Delta_{\text{evaluador}}$ es el artefacto del evaluador
canonico (siempre $\ge 0$) y $\gap_{\text{real}}$ es la deficiencia
real de la busqueda.

\subsection{Implementacion}

La implementacion se concentra en un solo modulo
(\texttt{scripts/\_re\_greedy\_20260529\_common.py}) que despacha por
los parametros CLI \texttt{--variant} y \texttt{--mh}. Un script
master de bash (\texttt{scripts/run\_all\_re\_greedy\_20260529.sh})
lanza las 25 grids (5 variantes en secuencia, 5 MHs en paralelo dentro
de cada variante).

\subsection{Smoke test}

Los smoke tests confirman que las 5 variantes funcionan en TS Simple
sobre gdb19 (1 corrida cada una):

\begin{itemize}
\item R1 strict\_greedy: 71 (gap 29.1\%).
\item R2 aos\_pm\_greedy: 65 (gap 18.2\%).
\item R3 pr\_greedy: \textbf{55 (BKS exacto)}.
\item R4 p\_inter\_pr\_greedy: \textbf{55 (BKS exacto)}.
\item R5 budget\_greedy: \textbf{55 (BKS exacto)}.
\end{itemize}

\subsection{Resultados completos}

El grid de 2875 corridas (5 variantes $\times$ 5 MHs $\times$ 23 instancias
$\times$ 5 repeticiones) se ejecuto con exito. La
Tabla~\ref{tab:re-greedy-gap} presenta el gap medio por MH y variante;
la Tabla~\ref{tab:re-greedy-bon} el \emph{Best-of-N} (mejor de 5
repeticiones por instancia, luego promediado); y la
Tabla~\ref{tab:re-greedy-bks} el numero de corridas que alcanzan el BKS
exacto (gap $< 0.01$\,\%).

\begin{table}[h]
\centering
\caption{Gap medio al BKS (\%) bajo evaluador greedy, por MH y variante.
         Sombreado: $< 1$\,\% (excelente), $1$--$5$\,\% (bueno),
         $5$--$10$\,\% (aceptable), $> 10$\,\% (deficiente).}
\label{tab:re-greedy-gap}
\small
\begin{tabular}{lrrrrrr}
\toprule
MH & R1 Strict & R2 AOS-PM & R3 PR & R4 p\_inter & R5 Budget & R6 WS* \\
\midrule
SA      & 15.11 & 15.02 & \textbf{0.90} & 2.13 & \textbf{0.57} & \textbf{0.00} \\
TS      & 15.48 & 14.52 & 4.48          & 5.22 & 1.08          & \textbf{0.00} \\
RTS     & 13.82 & 14.66 & 3.04          & 4.15 & 1.18          & \textbf{0.00} \\
ABC     & 14.98 & 14.82 & 9.75          & 10.74 & 8.30         & 1.32 \\
Cuckoo  & 14.81 & 14.20 & 8.20          & 8.97  & 4.25         & 7.08 \\
\midrule
\textbf{Media} & 14.84 & 14.64 & 5.27 & 6.24 & \textbf{3.08} & 1.68 \\
\bottomrule
\end{tabular}
\\[4pt]
\footnotesize *R6 warmstart-greedy: smoke test sobre 2 instancias (gdb19, kshs1).
\end{table}

\begin{table}[h]
\centering
\caption{Best-of-N gap (\%) — mejor de 5 repeticiones por instancia,
         luego promediado. Mide el techo de calidad de cada configuracion.}
\label{tab:re-greedy-bon}
\small
\begin{tabular}{lrrrrrr}
\toprule
MH & R1 Strict & R2 AOS-PM & R3 PR & R4 p\_inter & R5 Budget & R6 WS* \\
\midrule
SA      & 11.51 & 11.27 & \textbf{0.00} & 0.62 & \textbf{0.03} & \textbf{0.00} \\
TS      & 11.40 &  9.70 & 2.08          & 2.20 & 0.58          & \textbf{0.00} \\
RTS     & 10.07 & 10.80 & 1.20          & 1.61 & 0.48          & \textbf{0.00} \\
ABC     & 10.80 & 11.36 & 6.14          & 7.27 & 4.59          & 1.32 \\
Cuckoo  & 10.94 &  9.54 & 4.11          & 4.48 & 1.67          & 7.08 \\
\bottomrule
\end{tabular}
\end{table}

\begin{table}[h]
\centering
\caption{Corridas que alcanzan el BKS exacto (gap $< 0.01$\,\%) sobre
         $\approx 115$--$116$ corridas por celda (23 inst $\times$ 5 reps).}
\label{tab:re-greedy-bks}
\small
\begin{tabular}{lrrrrrr}
\toprule
MH & R1 & R2 & R3 & R4 & R5 & R6* \\
\midrule
SA      &  1 &  1 & \textbf{74} &  -- & \textbf{87} & 2/2 \\
TS      &  2 &  -- & 36         &  -- & 76          & 2/2 \\
RTS     &  0 &  -- & 38         &  -- & 72          & 2/2 \\
ABC     &  2 &  -- &  2         &  -- &  4          & 1/2 \\
Cuckoo  &  1 &  -- & 19         &  -- & 30          & 0/2 \\
\midrule
\textbf{Total} & 6 & -- & 169 & -- & 269 & 7/10 \\
\midrule
Instancias c/BKS & 4/23 & 5/23 & \textbf{23/23} & 18/23 & 22/23 & 2/2 \\
\bottomrule
\end{tabular}
\end{table}

El resultado mas destacable es que \textbf{PR (R3) alcanza el BKS exacto
en las 23 instancias del benchmark} (al menos en una de las 5
repeticiones), con SA como la MH mas efectiva (74 corridas en BKS,
64\,\% de tasa de exito individual). Budget (R5) aumenta la tasa de
SA al 75\,\% pero a un costo de tiempo $6.5\times$ mayor.

\paragraph{Analisis por instancia.}
La Tabla~\ref{tab:re-greedy-inst} muestra el gap medio sobre las 5 MHs
para las variantes R1, R3 y R5 en cada instancia. Las instancias mas
dificiles (\texttt{gdb12}, \texttt{gdb5}, \texttt{gdb21}) mantienen
gap $> 7$\,\% incluso en R5, lo que sugiere que el espacio de busqueda
en instancias con alta densidad de arcos y capacidad ajustada requiere
operadores mas agresivos o estrategias de reinicio global.

\begin{table}[h]
\centering
\caption{Gap medio al BKS (\%) por instancia — media sobre las 5 MHs.}
\label{tab:re-greedy-inst}
\small
\begin{tabular}{lrrrr}
\toprule
Instancia & BKS & R1 Strict & R3 PR & R5 Budget \\
\midrule
gdb1  &  316 & 12.65 & 3.47 & 1.89 \\
gdb2  &  339 & 23.46 & 6.51 & 4.38 \\
gdb3  &  275 & 19.75 & 5.63 & 3.00 \\
gdb4  &  287 & 19.18 & 4.43 & 2.66 \\
gdb5  &  377 & 28.08 & 5.19 & 3.78 \\
gdb6  &  298 & 11.50 & 4.13 & 2.97 \\
gdb7  &  325 &  6.26 & 2.31 & 0.91 \\
gdb10 &  275 & 15.21 & 8.51 & 4.65 \\
gdb12 &  458 & 37.02 & 13.24 & 8.59 \\
gdb13 &  536 & 14.96 & 11.13 & 7.47 \\
gdb14 &  100 & 17.12 & 3.76 & 2.64 \\
gdb15 &   58 &  9.52 & 1.24 & 0.83 \\
gdb16 &  127 & 10.83 & 6.49 & 3.34 \\
gdb17 &   91 &  2.11 & 0.53 & 0.35 \\
gdb19 &   55 & 17.20 & 3.22 & 1.21 \\
gdb20 &  121 & 14.98 & 7.97 & 4.20 \\
gdb21 &  156 & 20.77 & 9.95 & 7.41 \\
kshs1 & 14661 & 6.12 & 1.60 & 0.46 \\
kshs2 &  9863 & 7.17 & 0.59 & 0.30 \\
kshs3 &  9320 & 5.36 & 1.79 & 0.53 \\
kshs4 & 11498 & 14.33 & 9.28 & 4.69 \\
kshs5 & 10957 & 14.13 & 8.87 & 3.67 \\
kshs6 & 10197 & 13.53 & 1.47 & 1.21 \\
\midrule
\textbf{Media} & -- & 14.84 & 5.27 & 3.08 \\
\bottomrule
\end{tabular}
\end{table}

% =====================================================================
\section{Comparativa global del programa experimental}
\label{sec:comparativa-global}
% =====================================================================

\subsection{Tabla maestra de gap medio}

La Tabla~\ref{tab:comparativa-global} consolida el gap medio por MH y
experimento bajo el \textbf{evaluador greedy de orientacion} (corrida
R1--R5 completa). Los valores del primer ciclo (evaluador canonico)
eran artefactos de la funcion de costo y no se incluyen en la tabla
definitiva; se conservan en las secciones individuales de cada
experimento para documentar el proceso de descubrimiento.

\begin{table}[h]
\centering
\caption{Gap medio al BKS (\%) por MH y experimento (evaluador greedy,
         corrida completa R1--R5). Negrita: configuraciones $\leq 1$\,\%.}
\label{tab:comparativa-global}
\small
\begin{tabular}{lrrrrrr}
\toprule
MH & R1 Strict & R2 AOS-PM & R3 PR & R4 p\_inter & R5 Budget & R6 WS* \\
\midrule
SA      & 15.11 & 15.02 & \textbf{0.90} & 2.13 & \textbf{0.57} & \textbf{0.00} \\
TS      & 15.48 & 14.52 & 4.48          & 5.22 & \textbf{1.08} & \textbf{0.00} \\
RTS     & 13.82 & 14.66 & 3.04          & 4.15 & \textbf{1.18} & \textbf{0.00} \\
ABC     & 14.98 & 14.82 & 9.75          & 10.74 & 8.30         & 1.32 \\
Cuckoo  & 14.81 & 14.20 & 8.20          & 8.97  & 4.25         & 7.08 \\
\midrule
\textbf{Media} & 14.84 & 14.64 & 5.27 & 6.24 & \textbf{3.08} & 1.68 \\
\bottomrule
\end{tabular}
\\[4pt]
\footnotesize
*R6 warmstart-greedy: smoke test (2 instancias). Grid completo de 575 corridas pendiente.
\end{table}

\subsection{Tabla de tiempo medio por corrida}

La Tabla~\ref{tab:tiempo-medio} resume el costo computacional por
experimento (corrida R1--R5 bajo evaluador greedy). El evaluador greedy
agrega una comparacion de distancias por tarea pero no cambia la
complejidad asintotica; el incremento respecto de los tiempos originales
(evaluador canonico) es menor al 10\,\% en todos los casos. El
presupuesto $\times 25$ de R5 multiplica el tiempo por
$5$--$30\times$ segun la MH.

\begin{table}[h]
\centering
\caption{Tiempo medio por corrida (segundos) — evaluador greedy, corrida R1--R5.}
\label{tab:tiempo-medio}
\small
\begin{tabular}{lrrrrrr}
\toprule
MH & R1 Strict & R2 AOS & R3 PR & R4 p\_inter & R5 Budget & R6 WS* \\
\midrule
SA      &  28.1 &  32.4 &  40.6 &  37.1 & 260.5 &   6.4 \\
TS      &  14.4 &  17.7 &  21.3 &  19.6 & 354.8 &  25.7 \\
RTS     &  23.2 &  26.4 &  33.2 &  29.9 & 431.7 &  27.6 \\
ABC     &  11.1 &  14.5 &  29.7 &  26.9 & 197.6 &   5.2 \\
Cuckoo  &  13.8 &  17.7 &  20.0 &  18.2 & 216.9 &   6.9 \\
\bottomrule
\end{tabular}
\\[4pt]
\footnotesize *Tiempos de smoke test; no representativos del grid completo.
\end{table}

\subsection{Ranking final bajo evaluador greedy}

El ranking definitivo de configuraciones, ordenado por gap medio
sobre las 5 MHs y las 23 instancias del benchmark, es:

\begin{enumerate}
\item \textbf{R5 Budget}: 3.08\,\% — maximo presupuesto, mejor calidad global.
\item \textbf{R3 PR}: 5.27\,\% — relacion calidad/tiempo optima;
      SA alcanza BKS en el 64\,\% de corridas individuales.
\item R4 p\_inter\_pr: 6.24\,\% — inferior a R3 pese a mayor complejidad.
\item R1 Strict: 14.84\,\% — baseline corregido; AOS-PM (R2, 14.64\,\%)
      es estadisticamente equivalente.
\item R2 AOS-PM: 14.64\,\% — sin diferencia significativa respecto a Strict.
\end{enumerate}

\paragraph{Configuracion recomendada.}
Para el balance optimo entre complejidad de implementacion, tiempo de
computo y calidad de solucion, se recomienda \textbf{SA + Path Relinking
+ evaluador greedy (R3)}:

\begin{itemize}
\item Gap medio SA: 0.90\,\% (mediana: 0.00\,\%).
\item BKS exacto en 74/115 corridas (64\,\%) con $\sim 40$ segundos promedio.
\item Componentes adicionales sobre el baseline: dos monkey-patches,
      cero parametros nuevos que requieran tuning.
\item R5-Budget mejora 0.33\,pp sobre R3 en SA pero consume $6.5\times$
      mas tiempo y requiere ajuste de 6 parametros adicionales.
\end{itemize}

\paragraph{Comportamiento por MH.}
SA es la MH dominante en todos los experimentos con evaluador greedy.
TS y RTS siguen de cerca bajo presupuesto extendido (R5: 1.08\,\% y
1.18\,\%). ABC y Cuckoo no superan el 4\,\% en ninguna configuracion,
lo que sugiere que su mecanismo de exploracion de enjambre no se
beneficia del mismo modo del evaluador mas preciso ni del presupuesto
ampliado, posiblemente por la escasa diversidad entre agentes en
instancias de tamano pequeno--mediano.

% =====================================================================
\section{Conclusiones del capitulo}
\label{sec:conclusiones}
% =====================================================================

\subsection{Aprendizajes tecnicos}

\begin{enumerate}
\item \textbf{El evaluador es el componente mas critico.} Con el
  evaluador greedy de orientacion, el gap del baseline Strict baja de
  $\approx 31$\,\% a 14.84\,\% sin modificar ningun parametro de
  busqueda. Ninguno de los demas componentes experimentales (AOS, PR,
  p\_inter, budget) produce una mejora unitaria comparable.
\item \textbf{Path Relinking es la mejora mas eficiente en calidad/tiempo.}
  R3-PR reduce el gap en 9.57\,pp respecto a R1-Strict con solo
  $+12$\,s adicionales en SA. SA alcanza el BKS exacto en 23/23
  instancias y en el 64\,\% de corridas individuales.
\item \textbf{El presupuesto extendido es util pero costoso.} R5-Budget
  mejora 0.33\,pp sobre R3 en SA, pero consume $6.5\times$ mas tiempo
  (260\,s vs 40\,s). La relacion es favorable si el tiempo no es
  restriccion; desfavorable en escenarios de muchas instancias.
\item \textbf{AOS-PM no converge en el presupuesto disponible.} La
  diferencia entre R2 (14.64\,\%) y R1 (14.84\,\%) no es significativa.
  Las recompensas demasiado escasas ($\sim 1$\,\% de iteraciones mejoran)
  impiden que los pesos del AOS se especialicen.
\item \textbf{ABC y Cuckoo tienen un techo de calidad inferior.} En
  ninguna configuracion bajan del 4\,\%, frente al 0.57\,\% de SA.
  Su mecanismo de exploracion colectiva no compensa la perdida de
  especializacion local que logran las MHs de trayectoria.
\item \textbf{Path-Scanning produce soluciones de calidad con orientacion
  greedy.} La heuristica de Golden et al.~\cite{golden1983} alcanza gap $\approx 9$\,\%
  en las 23 instancias sin ninguna busqueda local, lo que la hace
  un punto de partida valioso para estudios futuros de ILS/LNS.
\end{enumerate}

\subsection{Aprendizajes metodologicos}

El aprendizaje metodologico mas importante del programa es el de la
Seccion~\ref{sec:evaluador-canonico}: una observacion empirica
aparentemente anecdota (que dos evaluadores diferian en $\sim 33$\,pp
sobre la misma solucion) revelo un artefacto en el evaluador rapido
del proyecto que estaba inflando todos los gaps de manera sistematica
y por igual en las cinco MHs.

Las lecciones metodologicas que se extraen:

\begin{enumerate}
\item \textbf{Validar el evaluador antes que la busqueda}: los gaps
  observados se interpretaron consistentemente como ``la busqueda no
  llega'' cuando la causa real era ``la medicion esta mal''. Una
  prueba ortogonal del evaluador con una construccion alternativa
  habria revelado el problema en el primer experimento.
\item \textbf{Atender a coincidencias sospechosas}: que las cinco
  MHs convergieran al mismo gap por instancia (sin importar su
  filosofia interna) era una senal fuerte de que el problema no
  estaba en las MHs. Esta senal se interpreto inicialmente como
  ``robustez'' cuando era ``firma de un artefacto comun''.
\item \textbf{El monkey-patching habilita comparaciones limpias}:
  poder activar/desactivar un componente sin tocar el codigo base de
  las MHs fue determinante para diagnosticar el artefacto del
  evaluador. La arquitectura por capas, aunque inicialmente disenada
  por comodidad, demostro su valor metodologico.
\end{enumerate}

\subsection{Lineas de trabajo abiertas}

\begin{enumerate}
\item \textbf{Grid completo de warmstart-greedy (R6).} La corrida
  actual cubre 2 instancias (smoke test). Ejecutar el grid de
  575 corridas (23 inst $\times$ 5 reps $\times$ 5 MH) permitira
  validar si la combinacion PS-warmstart + evaluador greedy + presupuesto
  extendido reduce el gap por debajo del 2\,\% en todo el benchmark.
\item \textbf{Analisis de instancias dificiles.} Las instancias
  \texttt{gdb12} (8.59\,\%), \texttt{gdb13} (7.47\,\%) y
  \texttt{gdb21} (7.41\,\%) mantienen gap elevado incluso en R5.
  Investigar si operadores mas agresivos (or-opt-3, GENIUS~\cite{gendreau1992}) o
  estrategias de reinicio global reducen ese piso.
\item \textbf{Comparacion greedy local vs optimo global de orientacion.}
  El evaluador greedy decide la orientacion de cada arco solo con
  informacion del nodo previo. El optimo global por ruta puede
  calcularse con programacion dinamica O(n). Cuantificar la diferencia
  en las 23 instancias y evaluar si justifica la complejidad adicional.
\item \textbf{Extension del benchmark.} Verificar la robustez de las
  conclusiones en instancias mayores (egl-s4, val9B) donde el espacio
  de busqueda es significativamente mayor.
\item \textbf{ILS sobre Path-Scanning.} Combinar PS con perturbacion
  sistematica (Iterated Local Search~\cite{lourenco2003}) para explotar el bajo gap inicial
  ($\approx 9$\,\%) del constructor heuristico sin la sobrecarga del
  ciclo de reheat de SA.
\end{enumerate}

% =====================================================================
\appendix
% =====================================================================

\section{Listado de modulos y scripts del programa}

\begin{table}[h]
\centering
\caption{Modulos del paquete metacarp asociados a cada experimento.}
\small
\begin{tabular}{ll}
\toprule
Experimento & Modulo \\
\midrule
strict\_intra\_inter & \texttt{metacarp/strict\_intra\_inter\_20260524.py} \\
lambda\_grid & (kwargs en wrappers) \\
aos\_pm & \texttt{metacarp/aos\_pm\_20260527.py} \\
path\_relinking & \texttt{metacarp/path\_relinking\_20260528.py} \\
p\_inter\_pr & \texttt{metacarp/p\_inter\_pr\_20260528.py} \\
budget & (kwargs + patch de $\lambda$) \\
warmstart\_greedy & \texttt{metacarp/path\_scanning.py} \\
                  & \texttt{metacarp/warmstart\_20260529.py} \\
                  & \texttt{metacarp/evaluador\_greedy\_20260529.py} \\
R1--R5 & \texttt{scripts/\_re\_greedy\_20260529\_common.py} \\
\bottomrule
\end{tabular}
\end{table}

\section{Instrucciones de reproducibilidad}

\subsection{Replicar un experimento individual}

\begin{lstlisting}[language=bash]
# Ejemplo: replicar p_inter_pr para SA
python scripts/run_sa_p_inter_pr_20260528.py --repeticiones 5

# Ejemplo: replicar warmstart_greedy para todas las MH en paralelo
for mh in sa tabu_simple tabu_reactiva abc_simple cuckoo; do
  python scripts/run_${mh}_warmstart_greedy_20260529.py \
    --repeticiones 5 \
    > logs/wsgreedy_${mh}.log 2>&1 &
done
wait
\end{lstlisting}

\subsection{Replicar el programa R1--R5 completo}

\begin{lstlisting}[language=bash]
# Master script: 25 grids = 5 variantes x 5 MH = 2875 corridas
bash scripts/run_all_re_greedy_20260529.sh
\end{lstlisting}

\section{Notas sobre el evaluador greedy y el optimo global de orientacion}

El evaluador greedy implementado decide la orientacion de cada arco en
funcion del nodo previo solamente. Esta decision es localmente optima
pero NO globalmente optima: en una ruta, la orientacion optima de un
arco puede depender no solo del previo sino tambien del siguiente.

El optimo global de orientacion para una ruta fija puede computarse
con programacion dinamica O(n) sobre dos estados por tarea (terminar
en $u$ o terminar en $v$). La diferencia tipica entre la version
greedy local y el optimo global es menor al 5\% en instancias de
tamano gdb/kshs; para benchmarks mas grandes podria ser relevante
considerarla.

